\documentclass{article}

% Language setting
% Replace `english' with e.g. `spanish' to change the document language
\usepackage[czech]{babel}
\usepackage{amsthm,thmtools,xcolor,amsmath,amssymb, mathtools}

% Set page size and margins
% Replace `letterpaper' with `a4paper' for UK/EU standard size
\usepackage[a4paper,top=2cm,bottom=2cm,left=3cm,right=3cm,marginparwidth=1.75cm]{geometry}

% Useful packages
\usepackage{amsmath}
\usepackage{enumerate}
\usepackage{graphicx}
\usepackage[colorlinks=true, allcolors=blue]{hyperref}
\graphicspath{images/}

%\usepackage{tikzit}
%\input{default.tikzstyles}

\title{DÚ Lineární algebra -- Sada 2}
\author{Jan Romanovský}

\begin{document}
\maketitle

\textbf{(2.1)} Napíšeme si skalární součiny pomocí matic a použijeme tvrzení 8.50, budeme hledat matici $A$:
$$\begin{pmatrix}
  x_1 & x_2
\end{pmatrix}A\begin{pmatrix}
  y_1 \\ y_2
\end{pmatrix}= \left\langle\begin{pmatrix}
  x_1 \\ x_2
\end{pmatrix}, \begin{pmatrix}
  y_1 \\ y_2
\end{pmatrix} \right\rangle \stackrel{8.50}{=}\left[\begin{pmatrix}
  x_1 \\ x_2
\end{pmatrix}\right]_B^T\left[\begin{pmatrix}
  y_1 \\ y_2
\end{pmatrix}\right]_B=\left([\operatorname{id}]_B^K\begin{pmatrix}
  x_1 \\ x_2
\end{pmatrix}\right)^T[\operatorname{id}]_B^K\begin{pmatrix}
  y_1 \\ y_2
\end{pmatrix}=$$ $$ = \begin{pmatrix}
  x_1 & x_2
\end{pmatrix} \left([\operatorname{id}]_B^K\right)^T[\operatorname{id}]_B^K\begin{pmatrix}
  y_1 \\ y_2
\end{pmatrix} = \begin{pmatrix}
  x_1 & x_2
\end{pmatrix} \left(\left([\operatorname{id}]_K^B\right)^{-1}\right)^T\left([\operatorname{id}]_K^B\right)^{-1}\begin{pmatrix}
  y_1 \\ y_2
\end{pmatrix}$$

Matici přechodu $[\operatorname{id}]_K^B$ můžeme přímo napsat podle bázových vektorů báze $B$. Odsud přímo vidíme vyjádření hledaného skalárního součinu maticově:
$$\begin{pmatrix}
  x_1 & x_2
\end{pmatrix}\left(\begin{pmatrix}
  1 & 3 \\ 2 & 7
\end{pmatrix}^{-1}\right)^T\begin{pmatrix}
  1 & 3 \\ 2 & 7
\end{pmatrix}^{-1}\begin{pmatrix}
  y_1 \\ y_2
\end{pmatrix}= \begin{pmatrix}
  x_1 & x_2
\end{pmatrix}\begin{pmatrix}
  7 & -2 \\ -3 & 1
\end{pmatrix}\begin{pmatrix}
  7 & -3 \\ -2 & 1
\end{pmatrix}\begin{pmatrix}
  y_1 \\ y_2
\end{pmatrix}=$$ $$=\begin{pmatrix}
  x_1 & x_2
\end{pmatrix}\begin{pmatrix}
  53 & -23 \\ -23 & 10
\end{pmatrix}\begin{pmatrix}
  y_1 \\ y_2
\end{pmatrix}$$
\end{document}
